\chapter{BackupEngine 与 Pipeline 接线}
\label{ch:engine-pipeline}

\section{Pipeline 内核选择}

\texttt{BackupEngine::RunBackup} 在 pipeline 模式下设置：
\begin{lstlisting}[caption={gtcdc\_kernel 注入 pipeline}]
if (CdcGtCdcEnabled()) {
  pipe_opts.gtcdc_kernel = GtCdcKernelForRepo(sb_);
}
\end{lstlisting}

\texttt{StreamingChunkCpuPipeline} 根据 \texttt{gtcdc\_kernel} 分发至 \texttt{Process2FGearChunkCut}（v6）或 \texttt{ProcessAnGearChunkCut}（v5）等。

\section{整文件 vs Streaming 路径}

\begin{table}[htbp]
\centering
\caption{v6 备份路径选择（与 FastCDC 规则对齐）}
\small
\begin{tabular}{@{}ll@{}}
\toprule
场景 & CDC 实现 \\
\midrule
Pipeline 大文件 sequential read & Streaming + \texttt{Process2FGearChunkCut} \\
Pipeline 小文件 / 内存映射 & \texttt{GtCdcSlice::ChunkCuts} \\
增量 \texttt{kIncremental} 大文件 & \texttt{ChunkFileFullGtCdc}（非 streaming） \\
\bottomrule
\end{tabular}
\end{table}

\section{Phase Stats 与 profile}

\texttt{EBBACKUP\_PIPELINE\_PROFILE=1} 时输出 \texttt{pipeline\_profile} 行；G-TCDC 额外细分：
\begin{itemize}[nosep]
  \item \texttt{stream\_cdc\_probes} / \texttt{gtcdc\_scan\_probes}
  \item \texttt{stream\_cdc\_ns}（计入 scan 次门禁分母）
\end{itemize}

\section{环境变量合约（v6 相关）}

\begin{longtable}{@{}p{4.5cm}p{9.5cm}@{}}
\toprule
\textbf{变量} & \textbf{语义} \\
\midrule
\endhead
\texttt{EBBACKUP\_CDC\_ALGO=gtcdc} & 启用 G-TCDC；必须与 G-TCDC 仓库配对 \\
\texttt{EBBACKUP\_CDC\_HYBRID=0} & bench 强制纯路径 \\
\texttt{EBBACKUP\_CDC\_FAST\_SLICE=0} & 禁用 FastSlice \\
\texttt{EBBACKUP\_PIPELINE\_PROFILE=1} & 启用 phase timing \\
\bottomrule
\caption{v6 bench 常用 env}
\end{longtable}
